\begin{abstract}
高效大语言模型（LLM）推理领域正在快速演进，机遇与挑战并存。
尽管该方向发展活跃，但长期缺少一个足够简明且可操作的统一分析框架，来帮助研究者系统理解各类 LLM 推理方法。
我们的综述不同于传统文献回顾：不仅总结研究现状，还提出了一个基于 Roofline 模型的系统化分析框架，用于刻画 LLM 推理中的关键瓶颈。
该框架可解释部署中的核心现实问题，例如：为何 LLM 往往是 memory-bound、其计算与内存需求有多大、如何选择合适硬件。
我们系统梳理了高效 LLM 推理的最新进展，覆盖模型压缩（如量化）、算法优化（如 speculative decoding）、系统与硬件层增强（如算子融合）等关键方向。
通过 Roofline 视角统一分析这些方法对内存访问与计算的影响，本综述既刻画研究全景，也为实际部署提供可落地的参考。
因此，本工作可作为新入门研究者的实践入口，也可帮助资深研究者深化对高效 LLM 部署问题的理解。
分析工具 \href{https://github.com/hahnyuan/LLM-Viewer}{LLM-Viewer} 已开源。
\end{abstract}
